\chapter{Estatística e probabilidade}\label{ch:g9:statproba}

Dados e acaso estão por toda parte: resultados esportivos, previsões do
tempo, jogos. Este capítulo ensina a resumir uma série de números pela
\omterm{def:g9:statproba:indicators}{média}, pela \omterm{def:g9:statproba:indicators}{mediana} e pela \omterm{def:g9:statproba:indicators}{amplitude}, e a calcular a \omterm{def:g9:statproba:probability}{probabilidade} de
experimentos aleatórios simples --- inclusive os de duas etapas,
tratados com diagramas de árvore. Os dois assuntos são desenvolvidos
muito mais longe no volume do ensino médio desta série.

\section{Resumir dados}

\begin{definition}[Média, mediana, amplitude]\label{def:g9:statproba:indicators}
Para uma série de $N$ números:
\begin{itemize}
  \item a \emph{média}\index{média} é a \omterm{def:g1:addition:def}{soma} de todos os valores
        dividida por $N$;
  \item a \emph{mediana}\index{mediana} é um valor que reparte a série
        ordenada em duas \omterm{def:g3:division:half}{metades} do mesmo tamanho: para $N$ \omterm{def:g3:numbers:evenodd}{ímpar}, o
        valor central; para $N$ par, o ponto médio dos dois valores
        centrais;
  \item a \emph{amplitude}\index{amplitude} é o maior valor menos o
        menor.
\end{itemize}
\end{definition}

\begin{example}\label{ex:g9:statproba:indicators}
Notas de um aluno: $8,\ 12,\ 9,\ 15,\ 11$.
\begin{enumerate}
  \item \omterm{def:g9:statproba:indicators}{Média}: $\dfrac{8 + 12 + 9 + 15 + 11}{5} = \dfrac{55}{5} = 11$.
  \item \omterm{def:g9:statproba:indicators}{Mediana}: ordene a série: $8, 9, 11, 12, 15$; o valor central (o
        3º) é $11$.
  \item \omterm{def:g9:statproba:indicators}{Amplitude}: $15 - 8 = 7$.
\end{enumerate}
Com uma sexta nota, $17$: ordenada, $8, 9, 11, 12, 15, 17$, a \omterm{def:g9:statproba:indicators}{mediana}
passa a ser o ponto médio de $11$ e $12$, isto é, $11.5$, e a \omterm{def:g9:statproba:indicators}{média}
passa a $\frac{72}{6} = 12$.
\end{example}

\begin{example}[Média ponderada]\label{ex:g9:statproba:weighted}
Os números de calçado vendidos num dia:
\begin{center}
\renewcommand{\arraystretch}{1.2}
\begin{tabular}{c|ccccc}
número & $38$ & $39$ & $40$ & $41$ & $42$ \\
\hline
quantidade & $3$ & $7$ & $6$ & $3$ & $1$ \\
\end{tabular}
\end{center}
O número médio é
\[
\frac{3 \times 38 + 7 \times 39 + 6 \times 40 + 3 \times 41 + 1 \times 42}
{3 + 7 + 6 + 3 + 1}
= \frac{114 + 273 + 240 + 123 + 42}{20}
= \frac{792}{20} = 39.6 .
\]
\end{example}

\begin{remark}
A \omterm{def:g9:statproba:indicators}{média} e a \omterm{def:g9:statproba:indicators}{mediana} podem diferir bastante. Na série
$10, 10, 10, 10, 60$, a \omterm{def:g9:statproba:indicators}{média} é $20$ (puxada para cima pelo valor
$60$), enquanto a \omterm{def:g9:statproba:indicators}{mediana} é $10$. Pergunte sempre qual indicador
representa melhor a situação.
\end{remark}

\section{Probabilidade}

\begin{definition}[Probabilidade]\label{def:g9:statproba:probability}
Um experimento aleatório tem vários \emph{resultados} possíveis; um
\emph{evento} é um conjunto de resultados. Cada resultado recebe uma
\emph{probabilidade}\index{probabilidade}: um número entre $0$ e $1$,
com todas as probabilidades dos resultados somando $1$; a probabilidade
$\P(A)$ de um evento $A$ é a \omterm{def:g1:addition:def}{soma} das probabilidades dos resultados
dele. Quando todos os $n$ resultados são igualmente prováveis,
\[
\P(A) = \frac{\text{número de resultados favoráveis}}{n}.
\]
\end{definition}

\begin{example}\label{ex:g9:statproba:spinner}
Uma roleta é dividida em $8$ setores iguais: $3$ vermelhos, $3$ azuis,
$2$ verdes. Cada setor tem \omterm{def:g9:statproba:probability}{probabilidade} $\frac18$, então
\[
\P(\text{vermelho}) = \frac38, \qquad
\P(\text{verde}) = \frac28 = \frac14, \qquad
\P(\text{não verde}) = 1 - \frac14 = \frac34 .
\]
\end{example}

\begin{proposition}[Complementar]\label{prop:g9:statproba:complement}
Para todo evento $A$, o \emph{evento contrário} $\bar A$ (``$A$ não
ocorre'') satisfaz
\[
\P(\bar A) = 1 - \P(A).
\]
\end{proposition}

\begin{proof}
Todo resultado está em exatamente um entre $A$ e $\bar A$, então
$\P(A) + \P(\bar A)$ \omterm{def:g1:addition:def}{soma} as \omterm{def:g9:statproba:probability}{probabilidades} de todos os resultados, o
que dá $1$.
\end{proof}

\section{Experimentos de duas etapas}

\begin{method}[Diagramas de árvore]\label{met:g9:statproba:tree}
Para um experimento em duas etapas:
\begin{enumerate}
  \item trace um galho por resultado da primeira etapa e depois
        continue cada galho com os resultados da segunda etapa,
        escrevendo cada \omterm{def:g9:statproba:probability}{probabilidade} no galho dela;
  \item multiplique as \omterm{def:g9:statproba:probability}{probabilidades} ao longo de um caminho para obter
        a \omterm{def:g9:statproba:probability}{probabilidade} desse caminho;
  \item some as \omterm{def:g9:statproba:probability}{probabilidades} de todos os caminhos que formam um
        evento.
\end{enumerate}
\end{method}

\begin{example}\label{ex:g9:statproba:tree}
Um saco tem $2$ fichas vermelhas e $3$ pretas. Tire uma ficha, anote a
cor, \emph{devolva-a} e tire de novo. Cada retirada dá vermelha com
\omterm{def:g9:statproba:probability}{probabilidade} $\frac25$ e preta com \omterm{def:g9:statproba:probability}{probabilidade} $\frac35$.

\begin{omfigure}
\begin{tikzpicture}[scale=0.95]
  \coordinate (root) at (0,0);
  \coordinate (R) at (2.2,1.0);
  \coordinate (B) at (2.2,-1.0);
  \coordinate (RR) at (4.7,1.55);
  \coordinate (RB) at (4.7,0.45);
  \coordinate (BR) at (4.7,-0.45);
  \coordinate (BB) at (4.7,-1.55);
  \draw (root) -- (R) node[midway, above, font=\small, sloped] {$\frac25$};
  \draw (root) -- (B) node[midway, below, font=\small, sloped] {$\frac35$};
  \draw (R) -- (RR) node[midway, above, font=\small, sloped] {$\frac25$};
  \draw (R) -- (RB) node[midway, below, font=\small, sloped] {$\frac35$};
  \draw (B) -- (BR) node[midway, above, font=\small, sloped] {$\frac25$};
  \draw (B) -- (BB) node[midway, below, font=\small, sloped] {$\frac35$};
  \node[omThm, font=\small, above left=-2pt] at (R) {V};
  \node[font=\small, below left=-2pt] at (B) {P};
  \node[right, font=\small] at (RR)
    {VV:\ $\frac25 \times \frac25 = \frac{4}{25}$};
  \node[right, font=\small] at (RB)
    {VP:\ $\frac25 \times \frac35 = \frac{6}{25}$};
  \node[right, font=\small] at (BR)
    {PV:\ $\frac35 \times \frac25 = \frac{6}{25}$};
  \node[right, font=\small] at (BB)
    {PP:\ $\frac35 \times \frac35 = \frac{9}{25}$};
\end{tikzpicture}

{\small A árvore das duas retiradas com reposição: multiplique ao longo
dos galhos e confira que os quatro caminhos somam $1$.}
\end{omfigure}

\omterm{def:g9:statproba:probability}{Probabilidade} de duas fichas da mesma cor:
$\frac{4}{25} + \frac{9}{25} = \frac{13}{25}$. \omterm{def:g9:statproba:probability}{Probabilidade} de pelo
menos uma vermelha:
$1 - \P(\text{PP}) = 1 - \frac{9}{25} = \frac{16}{25}$.
\end{example}

\begin{remark}[As frequências se aproximam das probabilidades]
Lançar um dado honesto $6000$ vezes dá \emph{cerca de} $1000$ seis ---
não exatamente. À medida que o número de repetições cresce, as
frequências observadas ficam cada vez mais próximas das \omterm{def:g9:statproba:probability}{probabilidades}:
é isso que faz da \omterm{def:g9:statproba:probability}{probabilidade} a ferramenta certa para modelar
experimentos repetidos (uma história desenvolvida no volume do ensino
médio desta série).
\end{remark}

\section{Exercícios}

\begin{exercise}[$\star$]\label{exo:g9:statproba:1}
Calcule a \omterm{def:g9:statproba:indicators}{média}, a \omterm{def:g9:statproba:indicators}{mediana} e a \omterm{def:g9:statproba:indicators}{amplitude} da série
$14,\ 9,\ 17,\ 9,\ 12,\ 11,\ 12$.
\end{exercise}

\begin{exercise}[$\star$]\label{exo:g9:statproba:2}
As temperaturas ao meio-dia ao longo de uma semana foram
$19,\ 21,\ 24,\ 18,\ 22,\ 25,\ 19$ (graus). Calcule a \omterm{def:g9:statproba:indicators}{média} (ao
décimo) e a \omterm{def:g9:statproba:indicators}{mediana}.
\end{exercise}

\begin{exercise}[$\star$]\label{exo:g9:statproba:3}
O número de gols marcados por um time em $10$ partidas:
\begin{center}
\renewcommand{\arraystretch}{1.2}
\begin{tabular}{c|cccc}
gols & $0$ & $1$ & $2$ & $3$ \\
\hline
partidas & $2$ & $4$ & $3$ & $1$ \\
\end{tabular}
\end{center}
Calcule o número médio de gols por partida e a \omterm{def:g9:statproba:indicators}{mediana}.
\end{exercise}

\begin{exercise}[$\star$]\label{exo:g9:statproba:4}
Um dado honesto é lançado uma vez. Calcule a \omterm{def:g9:statproba:probability}{probabilidade} de: ``sair
um $3$''; ``sair um \omterm{def:g3:numbers:evenodd}{número par}''; ``sair pelo menos $3$''; ``não sair
$6$''.
\end{exercise}

\begin{exercise}[$\star$]\label{exo:g9:statproba:5}
Uma caixa contém $12$ bolas: $5$ brancas, $4$ vermelhas e $3$ verdes;
uma bola é sorteada. Calcule $\P(\text{branca})$,
$\P(\text{vermelha ou verde})$ e $\P(\text{não branca})$. O que você
percebe sobre as duas últimas?
\end{exercise}

\begin{exercise}[$\star\star$]\label{exo:g9:statproba:6}
Uma moeda é lançada e depois um dado honesto é lançado. Trace a árvore
(ou conte os resultados) e calcule a \omterm{def:g9:statproba:probability}{probabilidade} de sair cara \emph{e}
um $6$; e depois a de sair cara \emph{ou} um $6$ (ou os dois).
\end{exercise}

\begin{exercise}[$\star\star$]\label{exo:g9:statproba:7}
No saco do \cref{ex:g9:statproba:tree} ($2$ vermelhas, $3$ pretas), as
duas retiradas agora são feitas \emph{sem} reposição. Trace a nova
árvore (atenção: as \omterm{def:g9:statproba:probability}{probabilidades} do segundo nível mudam) e calcule a
\omterm{def:g9:statproba:probability}{probabilidade} de tirar duas fichas pretas e depois a de tirar duas
fichas de cores diferentes.
\end{exercise}

\begin{exercise}[$\star\star$]\label{exo:g9:statproba:8}
Depois de $9$ partidas, uma jogadora de basquete tem \omterm{def:g9:statproba:indicators}{média} de $14$
pontos por partida.
\begin{enumerate}
  \item Quantos pontos ela marcou no total?
  \item Quantos pontos ela precisa marcar na $10^{\text{a}}$ partida
        para ficar com \omterm{def:g9:statproba:indicators}{média} $15$?
\end{enumerate}
\end{exercise}

\begin{exercise}[$\star\star\star$]\label{exo:g9:statproba:9}
Um jogo: lance dois dados honestos; você ganha se os dois resultados
forem iguais (``uma dupla'').
\begin{enumerate}
  \item Calcule a \omterm{def:g9:statproba:probability}{probabilidade} de ganhar.
  \item Você joga duas vezes (partidas independentes). Usando uma
        árvore, calcule a \omterm{def:g9:statproba:probability}{probabilidade} de ganhar pelo menos uma vez.
\end{enumerate}
\end{exercise}

\section{Problema: as apostas ruinosas do Cavaleiro}

\begin{problem}[{Problema de fim de semana --- a disputa de jogo de
1654 que criou a teoria das probabilidades e a coincidência de
aniversários que engana toda sala de aula}]
\label{pb:g9:statproba:1}
Em 1654, um nobre jogador inveterado, o Cavaleiro de Méré,
enriqueceu com uma aposta de dados e começou a perder em outra que
acreditava ser equivalente. Perplexo, ele perguntou ao matemático
Blaise Pascal; Pascal escreveu a Pierre de Fermat; e a troca de cartas
entre eles fundou a teoria das \omterm{def:g9:statproba:probability}{probabilidades}. Este problema encena de
novo o caso inteiro com as ferramentas deste capítulo --- resultados
igualmente prováveis, árvores e a regra do complementar
(\cref{prop:g9:statproba:complement}) --- e termina com a mais
contraintuitiva de todas as coincidências.

\medskip
\noindent\textbf{Parte I --- Dados, contados com honestidade.}
\begin{enumerate}
  \item Lance um dado honesto duas vezes. Usando uma árvore (ou uma
        tabela $6 \times 6$), calcule a \omterm{def:g9:statproba:probability}{probabilidade} de \emph{nenhum}
        seis nos dois lançamentos e deduza a \omterm{def:g9:statproba:probability}{probabilidade} de pelo
        menos um seis.
  \item Um colega argumenta: ``uma chance em seis por lançamento, duas
        vezes, então $\frac16 + \frac16 = \frac13$''. Compare com a
        questão 1 e aponte os resultados exatos que a \omterm{def:g1:addition:def}{soma} dele conta
        duas vezes.
  \item Lançando dois dados e somando: calcule a \omterm{def:g9:statproba:probability}{probabilidade} de uma
        \omterm{def:g1:addition:def}{soma} $7$ e de uma \omterm{def:g1:addition:def}{soma} $2$. Por que o $7$ é o preferido dos
        jogadores?
  \item Disputa antiga: com dois dados, uma \omterm{def:g1:addition:def}{soma} $9$ ou uma \omterm{def:g1:addition:def}{soma} $10$ é
        mais provável? Conte os resultados e resolva a questão.
  \item Generalize a questão 1: explique por que a \omterm{def:g9:statproba:probability}{probabilidade} de
        pelo menos um seis em $n$ lançamentos é
        \[
        1 - \left(\frac56\right)^{\!n} ,
        \]
        já que os galhos sem seis da árvore se multiplicam de nível em
        nível.
\end{enumerate}

\medskip
\noindent\textbf{Parte II --- As duas apostas do Cavaleiro.}
\begin{enumerate}[resume]
  \item Primeira aposta, paga a um para um: \emph{pelo menos um seis em
        quatro lançamentos de um dado}. Calcule
        $\left(\frac56\right)^4$ como \omterm{def:g9:fractions:fraction}{fração} e como decimal, e a
        \omterm{def:g9:statproba:probability}{probabilidade} de vitória do Cavaleiro. A aposta era boa?
  \item O raciocínio do próprio Cavaleiro era: ``quatro lançamentos a
        $\frac16$ cada: $\frac46$ de chance''. Ele deu quase a resposta
        certa neste caso --- mas leve-o até $7$ lançamentos: que
        absurdo ele produz, e que erro da questão 2 ele repete?
  \item Segunda aposta, que ele acreditava equivalente por
        \omterm{def:g9:linfunc:linear}{proporcionalidade} (``$24$ lançamentos a $\frac1{36}$ cada: de
        novo $\frac{24}{36} = \frac46$''): \emph{pelo menos um seis
        duplo em $24$ lançamentos de dois dados}. Calcule a
        \omterm{def:g9:statproba:probability}{probabilidade} verdadeira de vitória,
        $1 - \left(\frac{35}{36}\right)^{24}$ (com a calculadora). Por
        que o Cavaleiro foi lentamente à ruína?
  \item Quantos lançamentos de dois dados teriam restaurado a vantagem
        dele? Teste $n = 25$ na calculadora e conclua.
  \item Enuncie a moral em duas frases: o que há de errado, em geral,
        em multiplicar uma \omterm{def:g9:statproba:probability}{probabilidade} pelo número de tentativas ---
        e que ferramenta correta substitui essa tentação?
\end{enumerate}

\medskip
\noindent\textbf{Parte III --- \omterm{def:g9:statproba:indicators}{Médias}, \omterm{def:g9:statproba:indicators}{medianas}, aniversários.}
\begin{enumerate}[resume]
  \item Uma empresa pequena paga $2\,000$ euros por mês a nove
        funcionários e $20\,000$ ao diretor. Calcule o salário médio e o
        salário mediano (\cref{def:g9:statproba:indicators}). Que número
        o anúncio de emprego deveria citar, honestamente?
  \item Construa uma lista de cinco notas de prova com \omterm{def:g9:statproba:indicators}{média} $12$ e
        \omterm{def:g9:statproba:indicators}{mediana} $15$. O que o par (\omterm{def:g9:statproba:indicators}{média} abaixo da \omterm{def:g9:statproba:indicators}{mediana}) diz sobre o
        formato das notas?
  \item O problema dos aniversários: numa turma de $23$ alunos, qual é
        a \omterm{def:g9:statproba:probability}{probabilidade} de pelo menos dois fazerem aniversário no mesmo
        dia? Monte o complementar (os $23$ aniversários todos
        diferentes):
        \[
        \frac{365}{365} \times \frac{364}{365} \times
        \frac{363}{365} \times \dots \times \frac{343}{365} .
        \]
        Calcule o \omterm{def:g2:mult:def}{produto} dos três primeiros fatores, descreva como os
        fatores evoluem e --- sabendo que o \omterm{def:g2:mult:def}{produto} completo é cerca de
        $0.493$ --- responda à pergunta. A maioria das pessoas espera
        ``muito improvável'': o que a matemática diz?
  \item A intuição consertada: quantos \emph{pares} de alunos podem
        compartilhar um aniversário numa turma de $23$ (a contagem de
        apertos de mão do \cref{pb:g6:wholes:1})? Explique numa frase
        por que é esse número, e não o $23$ em si, que produz a
        surpresa.
  \item Final: descreva um experimento para conferir um dos dois
        resultados --- a aposta de Méré com um dado de verdade ou o
        problema dos aniversários em várias turmas --- e enuncie com
        cuidado o que você espera das frequências observadas à medida
        que o número de repetições cresce. (Essa expectativa tem nome,
        lei dos grandes números; a demonstração dela espera no volume do
        ensino médio.)
\end{enumerate}
\end{problem}
